\documentclass[11pt]{article}

% --- Compile with: latexmk -xelatex proposal.tex ---
\usepackage{fontspec}          % XeLaTeX: native UTF-8, Vietnamese-capable
\setmainfont{Latin Modern Roman}
\usepackage{geometry}
\geometry{a4paper, margin=1in}
\usepackage{microtype}
\usepackage{booktabs}
\usepackage{tabularx}
\usepackage{array}
\usepackage{enumitem}
\setlist{nosep, leftmargin=1.4em}
\usepackage{xcolor}
\usepackage[colorlinks=true, linkcolor=black, citecolor=blue!55!black,
            urlcolor=blue!55!black]{hyperref}
\usepackage[backend=biber, style=numeric-comp, sorting=none, url=true,
            doi=false, isbn=false]{biblatex}
\addbibresource{references.bib}
\usepackage{titlesec}
\titleformat{\section}{\large\bfseries}{\thesection}{0.6em}{}
\titleformat{\subsection}{\normalsize\bfseries}{\thesubsection}{0.6em}{}
\titlespacing*{\section}{0pt}{1.1em}{0.5em}
\titlespacing*{\subsection}{0pt}{0.8em}{0.35em}

\newcolumntype{Y}{>{\raggedright\arraybackslash}X}

\title{\vspace{-2.2em}\textbf{Content Disarm and Reconstruction as an
Adversarial Defense for Vision-Language Models:\\
An Empirical Audit of Deployed CDR Tools}}
\author{Nông Văn Tình}
\date{Advanced Computer Vision (ACV), Project Proposal \\ \today}

\begin{document}
\maketitle
\vspace{-1.2em}

% ======================================================================
\section{Project Introduction}

\subsection{Background / Scope}
Vision-Language Models (VLMs) such as LLaVA~\cite{liu2023llava} couple a
CLIP-style Vision Transformer~\cite{radford2021clip} to a language model. This
visual channel is vulnerable to \emph{pixel-space adversarial perturbations}:
a small, bounded change to the image, optimized end-to-end through the model,
can steer the generated text~\cite{qi2024visualaaai,univperturb2024}. A common,
model-agnostic defense family is \emph{input transformation}: re-encoding or
reconstructing the image before it reaches the encoder~\cite{guo2018countering}.

Independently, the security industry deploys \emph{Content Disarm and
Reconstruction (CDR)}: parse a file, discard non-conforming or active structure
(metadata, embedded objects, anomalous chunks), and rebuild a clean file from
validated content. Real, reproducible open-source CDR tools already exist.
\emph{ICDR}~\cite{icdr,openimagecdr} has a published pipeline of
\texttt{transCode} (re-encode, strip all metadata) $\rightarrow$ \texttt{resize}
(97\% compress/restore) $\rightarrow$ \texttt{AdvFilter} (Gaussian blur then
sharpen). \emph{Dangerzone}~\cite{dangerzone} renders a document to
raw RGB pixels and rebuilds a clean file, discarding \emph{all} non-pixel
structure. This project evaluates these existing implementations.

\textbf{Scope.} This is a computer-vision robustness study: we audit whether
deployed image-CDR tools, applied as a pre-encoder defense, reduce the success of
targeted adversarial images against a VLM while preserving legitimate visual
content, and whether they do better than trivial JPEG re-encoding. The methods
are core CV: adversarial optimization through a ViT-based encoder, frequency-domain
analysis of what reconstruction removes, image-fidelity measurement, and
vision-encoder robustness evaluation.

\subsection{Problem Statement / Motivation}
A named security framework, SAF-M-49~\cite{safm49}, \emph{recommends} multimedia
CDR as a mitigation against multimodal prompt injection, but it provides no
empirical validation (it lists efficacy and latency measurement as future work).
Meanwhile, existing CDR-for-AI research sanitizes \emph{model
weight files}~\cite{aimodelcdr,ztaicdr}, not input images, and image-CDR research
targets \emph{malware / steganography removal}~\cite{openimagecdr}, not
adversarial robustness. Purpose-built adversarial \emph{purification} for VLMs
does exist, but it is bespoke machine learning (diffusion or latent-space
purifiers)~\cite{nie2022diffpure,clipure2025,codefend2025}, not the deterministic
security tools practitioners actually deploy. This leaves a concrete, testable
question: \emph{does a deployed CDR tool, recommended for exactly this
purpose, actually blunt adversarial images on a VLM, and is it worth more than a
one-line JPEG call?} My professional background in CDR motivates a faithful setup
and a mechanistic reading of \emph{which} disarming operation carries any effect.

% ======================================================================
\section{Related Work / Literature Survey}
We separate seven categories to avoid conflation.
\textbf{(1)~CDR on model files/weights}~\cite{aimodelcdr,ztaicdr}: different
surface. \textbf{(2)~CDR on input images}~\cite{icdr,openimagecdr,dangerzone}:
real tools, but evaluated only for malware/steganography, never adversarial
robustness. \textbf{(3)~Conventional transform
defenses}~\cite{guo2018countering,shin2017jpegresistant}: single-operation
(mostly JPEG), classifier-era. \textbf{(4)~Purification defenses for
VLMs}~\cite{nie2022diffpure,clipure2025,codefend2025}: numerous but bespoke ML,
not deployed security tools. \textbf{(5)~Adversarial attacks on
VLMs}~\cite{qi2024visualaaai,univperturb2024}: supply our attack methodology.
\textbf{(6)~Typographic / prompt injection}~\cite{scenetap2024,jailbreaksurvey2024}:
out of MVP scope (see \S5). \textbf{(7)~Robust-evaluation
methodology}~\cite{athalye2018obfuscated,tramer2020adaptive,carlini2019evaluating}:
governs our adaptive-attack and statistics choices.

\noindent\textbf{Positioned gap (narrow, supported).} We are not aware of any
empirical evaluation of \emph{deployed image-CDR tools} (ICDR, Dangerzone) as a
pre-encoder adversarial defense for VLMs, nor of a test of whether they provide
robustness \emph{beyond plain JPEG}. We make no ``first/novel'' claim about CDR as
a concept: SAF-M-49 already recommends it, and the contribution is the missing audit.

% ======================================================================
\section{Project Content Outline}

\subsection{Core Idea}
Insert a deployed CDR tool as a pre-encoder stage,
\texttt{image $\rightarrow$ CDR $\rightarrow$ VLM}, and measure targeted attack
success and clean-task utility against two references: no defense, and a JPEG
quality sweep that traces the security--utility frontier. Then ask whether CDR
lies beyond that frontier, and \emph{why}.

\subsection{Research Questions}
\begin{itemize}
\item \textbf{RQ1.} Does a deployed CDR tool reduce the targeted attack success
      rate of adversarial images on a VLM while preserving clean VQA accuracy?
\item \textbf{RQ2.} Does CDR provide robustness \emph{beyond} the JPEG
      quality-sweep frontier at comparable utility cost?
\item \textbf{RQ3 (mechanism).} Is any robustness explained by removal of
      high-frequency perturbation energy, i.e.\ does the causal chain
      \emph{CDR operation $\to$ perturbation-energy removal $\to$ attack failure}
      hold?
\item \textbf{RQ4 (extension).} Can a CDR-aware adaptive attacker recover attack
      success, indicating the non-adaptive result overstates robustness?
\end{itemize}

\subsection{Threat Model}
\begin{itemize}
\item \textbf{Attacker.} White-box on the target VLM (following
      Qi et al.~\cite{qi2024visualaaai}); may perturb the input image only, under
      an $L_\infty$ budget $\varepsilon$; optimizes teacher-forcing cross-entropy
      from raw pixels to LM logits toward a predetermined target answer.
\item \textbf{Knowledge (two settings).} \emph{Defense-unaware} (MVP: attack
      then sanitize) and \emph{defense-aware} (Extension: attack \emph{through}
      the CDR pipeline).
\item \textbf{Success criterion (single, primary).} \emph{Targeted VQA-answer
      flip}: for a controlled single-object question with known ground truth, the
      model outputs the attacker's predetermined incorrect label. No
      harmful-content generation.
\item \textbf{Defender.} Applies one sanitizer before the encoder; compared with
      no defense and with JPEG.
\end{itemize}

\subsection{Dataset}
200 images from the ImageNet-1k validation
set~\cite{deng2009imagenet}: 20 classes $\times$ 10 images (fixed seed list,
IDs logged). Each image gets one controlled prompt (``What is the main object?
Answer with a single label.'') with the ImageNet class as ground truth and one
\emph{predetermined} target label (fixed class-to-target map). Images are resized
to the model's native input (336\,px for LLaVA). The \emph{same} images pass
through \emph{every} defense (paired design). Attacks are optimized per image on
the test image itself, which is standard for per-example adversarial attacks.
Because no model is trained, there is no train/test leakage, only per-image
white-box optimization. Clean-utility is measured on the unattacked images under
each defense.

\subsection{Models}
\textbf{Primary: LLaVA-1.5-7B} (the Qi-et-al.~lineage, maximizing comparability
and reproducibility). No CLIP surrogate is used; the attack runs end-to-end on the
target model. \textbf{Compute budget:} a single 24\,GB GPU; white-box PGD through
7B parameters uses fp16 and gradient checkpointing. \textbf{Fallback} for
$<16$\,GB: Qwen2-VL-2B~\cite{wang2024qwen2vl} as primary. A second model is used
\emph{only} in the extension to test generalization, not in the MVP.

\subsection{Experimental Design}
\textbf{Defenses (minimal, principled):}
\texttt{D0} none; \texttt{D1--D3} JPEG at quality $\{90,75,50\}$ (traces the
frontier); \texttt{D4} ICDR (\texttt{final\_fun}: transCode$\to$resize$\to$AdvFilter)
\cite{icdr}; \texttt{D5} Dangerzone (full render-and-rebuild)~\cite{dangerzone}.
\textbf{Ablation (conditional).} \emph{Only if} ICDR (\texttt{D4}) differs
significantly from the JPEG frontier do we decompose its published operations
(transCode / resize / AdvFilter / the alternate LSB-clean and gamma steps) to
localize the effect. This keeps the matrix small unless the data justify
expansion.
\textbf{Comparison protocol.} Plot every defense in
(clean-accuracy-retained, robust-accuracy-under-attack) space; JPEG
$\{90,75,50\}$ traces the frontier; CDR tools appear as points. RQ2 asks whether
CDR lies \emph{on} or \emph{beyond} that frontier. Fidelity metrics are reported
as descriptive covariates, not as a matching axis.

\subsection{Step-by-Step To-Do Plan}
\begin{enumerate}
\item Load LLaVA-1.5-7B; wrap it as a differentiable pixel$\to$logits function.
\item Integrate deployed defenses: JPEG (Pillow), ICDR~\cite{icdr},
      Dangerzone~\cite{dangerzone}; verify each runs deterministically (log
      versions/configs).
\item Build the 200-image ImageNet VQA benchmark + fixed target-label map.
\item Implement the targeted PGD attack (teacher forcing to target answer),
      $\varepsilon\in\{4,8,16\}/255$; validate attack strength on \texttt{D0}.
\item Run the paired ladder: attack $\times$ \{\texttt{D0--D5}\} on the primary
      model. \textbf{$\leftarrow$ MVP result.}
\item Mechanistic analysis (RQ3): residual perturbation energy and its frequency
      spectrum before/after each defense.
\item Statistics: McNemar per defense pair, bootstrap CIs, Holm correction.
\item \textbf{Extension:} defense-aware adaptive attack (\S3.6, \S6); optional
      second model for generalization.
\item Aggregate, plot the frontier, build the demo, write up.
\end{enumerate}

% ======================================================================
\section{Results}

\subsection{Expected Outcomes}
A security--utility frontier that places deployed CDR tools relative to JPEG and
to no defense, on one VLM with a targeted white-box attack; a mechanistic test of
\emph{why} any effect occurs; an optional adaptive-attack check; and a
reproducible codebase and demo.

\subsection{Evaluation Metrics (each tied to one RQ)}
\begin{tabularx}{\textwidth}{@{}l Y@{}}
\toprule
\textbf{Metric} & \textbf{Research question it answers} \\
\midrule
Targeted ASR ($\pm$95\% bootstrap CI) & RQ1, RQ2, RQ4: does the defense stop the attack? \\
Clean VQA accuracy under each defense & RQ1: does the defense preserve utility? \\
Security--utility frontier plot & RQ2: does CDR beat JPEG at equal utility cost? \\
Residual perturbation energy + FFT spectrum & RQ3: does robustness come from removing high-freq perturbation? \\
PSNR / SSIM / LPIPS (descriptive) & context for the frontier (not a matching axis) \\
CDR/JPEG latency & practical cost of deployment \\
\bottomrule
\end{tabularx}

\subsection{Experimental Evaluation}
Paired design (identical images through all defenses) $\Rightarrow$
\textbf{McNemar's test} for each pairwise ASR comparison; \textbf{95\% bootstrap
CIs} on ASR and clean accuracy; \textbf{Holm--Bonferroni} across the small set of
comparisons; effect size reported as the ASR difference. Attack validity is
checked per Carlini et al.~\cite{carlini2019evaluating} (attack must reach high
ASR on \texttt{D0} before any defense claim).

\subsection{Visual Demo}
Panels of (clean $\to$ correct answer), (adversarial $\to$ target answer),
(adversarial $\to$ CDR $\to$ recovered or not), each annotated with its FFT
spectrum and the residual-energy readout; plus the frontier plot with the CDR
points marked.

\subsection{Hypotheses and Falsification}
\begin{itemize}
\item \textbf{H1:} CDR reduces targeted ASR vs.\ \texttt{D0}.
      \emph{Falsified if} ASR under \texttt{D4/D5} is within the CI of \texttt{D0}.
\item \textbf{H2:} CDR lies beyond the JPEG frontier for $\geq 1$ operating point.
      \emph{Falsified if} CDR points sit on/inside the JPEG curve.
\item \textbf{H3:} Robustness tracks high-frequency perturbation removal.
      \emph{Falsified if} defenses that remove comparable high-freq energy give
      divergent ASR, or vice versa.
\item \textbf{H4:} Aggressive CDR (Dangerzone) costs measurable clean accuracy.
      \emph{Falsified if} clean accuracy under \texttt{D5} is within \texttt{D0}'s CI.
\item \textbf{H5:} A defense-aware attack recovers substantial ASR.
      \emph{Falsified if} adaptive ASR stays near non-adaptive defended ASR.
\end{itemize}

% ======================================================================
\section{Scope and Limitations}
This project is an audit, not a new defense method. The MVP studies
\emph{one} attack family (pixel-space perturbation), \emph{one} primary VLM, and
\emph{deployed} CDR tools. Commercial CDR (e.g.\ MetaDefender, Votiro) is
proprietary and not bit-reproducible, so it is out of scope.
Success is a targeted-VQA proxy for output manipulation, not harmful-content
generation. Attacks are per-image white-box; results bound the white-box case,
the strongest for the defender to survive. Typographic and visual
prompt-injection is deliberately excluded from the MVP and left as an optional
contrast experiment. Rendered text is legitimate visible content, so CDR is
expected \emph{not} to remove it, and Dangerzone's OCR layer could even
re-inject it. We would only spot-check this, not build the project on it.

% ======================================================================
\section{Risks and Fallbacks}
\begin{tabularx}{\textwidth}{@{}Y Y@{}}
\toprule
\textbf{Risk} & \textbf{Fallback} \\
\midrule
White-box PGD through 7B exceeds GPU memory &
  Gradient checkpointing + fp16; else Qwen2-VL-2B as primary. \\
ICDR/Dangerzone hard to script in the pipeline &
  Run them as external processes on saved images (batch), then load; both are
  open-source and file-based. \\
CDR indistinguishable from JPEG &
  That \emph{is} the RQ2 answer (H2 falsified), a useful negative result for
  practitioners. \\
Adaptive attack (BPDA) too heavy through 7B $+$ CDR &
  Extension only; simpler fallback is a defense-aware transfer attack through a
  differentiable CDR approximation (differentiable-JPEG + straight-through for
  non-differentiable steps). \\
Attack too weak on \texttt{D0} &
  Increase $\varepsilon$/iterations until high ASR on \texttt{D0} before any
  defense claim~\cite{carlini2019evaluating}. \\
Time runs out &
  MVP (RQ1+RQ2, one model, one attack, 200 images) stands alone. \\
\bottomrule
\end{tabularx}

% ======================================================================
\section{Timeline}
\begin{tabularx}{\textwidth}{@{}l Y@{}}
\toprule
\textbf{Week} & \textbf{Deliverable} \\
\midrule
W1 & Differentiable VLM wrapper; integrate JPEG/ICDR/Dangerzone; benchmark built. \\
W2 & Targeted PGD implemented and validated on \texttt{D0}; paired ladder run. \textbf{$\leftarrow$ MVP result.} \\
W3 & Mechanistic analysis (residual energy, FFT) + statistics (McNemar, CIs, Holm) + frontier plot. \\
W4 & \emph{Extension (optional):} adaptive attack and/or second-model generalization; demo; write-up. \\
\bottomrule
\end{tabularx}

% ======================================================================
\section*{Appendix A: Conservative Novelty Statement}
\addcontentsline{toc}{section}{Appendix A: Novelty Statement}
Image Content Disarm and Reconstruction is an existing, deployed
file-sanitization technique, with open-source implementations such as
ICDR~\cite{icdr,openimagecdr} and Dangerzone~\cite{dangerzone}, and security
guidance already recommends multimedia CDR against multimodal prompt
injection~\cite{safm49}. However, that recommendation is unvalidated, existing
CDR-for-AI research sanitizes model weight files rather than model
inputs~\cite{aimodelcdr,ztaicdr}, and purpose-built adversarial purification for
VLMs is bespoke machine learning rather than a deployed security
tool~\cite{nie2022diffpure,clipure2025}. We therefore do not claim CDR is a novel
idea. Our contribution is an empirical audit: we are not aware of prior work that
evaluates deployed image-CDR tools as a pre-encoder adversarial defense for
vision-language models, or that tests whether they provide robustness beyond
plain JPEG re-encoding, with a controlled targeted white-box attack, a
security--utility frontier comparison, a mechanistic frequency-domain analysis,
and an adaptive-attack robustness check.

% ======================================================================
\section*{Appendix B: Experimental Matrix}
\addcontentsline{toc}{section}{Appendix B: Experimental Matrix}
\begin{tabularx}{\textwidth}{@{}Y l Y@{}}
\toprule
\textbf{Attack} & \textbf{Defenses / Model} & \textbf{Primary metrics} \\
\midrule
Pixel-space targeted PGD ($L_\infty$, $\varepsilon\in\{4,8,16\}/255$),
white-box on the VLM~\cite{qi2024visualaaai}
  & \texttt{D0} none; \texttt{D1--D3} JPEG$\{90,75,50\}$; \texttt{D4} ICDR;
    \texttt{D5} Dangerzone. Model: LLaVA-1.5-7B.
  & Targeted ASR ($\pm$CI); clean VQA accuracy; residual perturbation energy +
    FFT; PSNR/SSIM/LPIPS; latency \\
\addlinespace
Adaptive (defense-aware) \emph{[Extension]}
  & \texttt{D4/D5}; LLaVA-1.5-7B ($+$ Qwen2-VL-2B for generalization)
  & Adaptive ASR vs.\ non-adaptive defended ASR \\
\bottomrule
\end{tabularx}

% ======================================================================
\section*{Appendix C: MVP vs.\ Extension}
\addcontentsline{toc}{section}{Appendix C: MVP vs. Extension}
\textbf{MVP (must implement; a complete result on its own):} differentiable
LLaVA-1.5-7B wrapper; integrated deployed defenses (JPEG sweep, ICDR,
Dangerzone); 200-image ImageNet VQA benchmark; targeted white-box PGD validated
on \texttt{D0}; paired ladder \texttt{D0--D5}; security--utility frontier;
mechanistic residual-energy/FFT analysis; paired statistics; tables, frontier
plot, and demo. Answers RQ1--RQ3 and remains interesting even if CDR $=$ JPEG.

\smallskip
\noindent\textbf{Extension (optional):} defense-aware adaptive
attack~\cite{athalye2018obfuscated} testing for a false sense of security (RQ4);
generalization to a second VLM; a small typographic-injection contrast check.

\printbibliography[title={References}]

\end{document}
